% NichidaiMasterExam_R8_1st_3
\begin{tcolorbox} %R8_1st_3
	$\fbox{3}$ 次の問いに答えよ。
\begin{enumerate}
	\item 定積分$\int_0^{2\pi} \sin^2{\theta}d\theta$を求めよ。
	\item 変数変換$r= \sin{u}$を用いて、定積分$\int_0^1 r^2\sqrt{1- r^2}dr$を求めなさい。
	\item 変数変換$\begin{cases}x= r\cos{\theta},\\ y= 2r\sin{\theta}\end{cases}$に対して、ヤコビアン$\frac{\partial(x, y)}{\partial(r, \theta)}$を求めなさい。
	\item $D= \{(x, y);\ 4x^2+ y^2\le 4\}$に対して、$\iint_D (x^2+ y^2)dxdy$を求めなさい。
	\item 変数変換$\begin{cases}x= r\sin{\theta}\cos{\phi},\\ y= r\sin{\theta}\sin{\phi},\\ z= r\cos{\theta} \end{cases}$に対して、ヤコビアン$\frac{\partial(x, y, z)}{\partial(r, \theta, \phi)}$を求めなさい。
	\item $E= \{(x, y, z);\ x^2+ y^2+ z^2\le 1\}$に対して、$\iiint_E 2\sqrt{1- x^2- y^2- z^2}dxdydz$を求めなさい。
\end{enumerate}
\end{tcolorbox}

\begin{souanswer} %R8_1st_2_ans
\begin{enumerate}
    \item 倍角公式より
\begin{align*}
    \int_0^{2\pi} \sin^2{\theta}
    &= \int_0^{2\pi} \frac{1- \cos{2\theta}}{2}d\theta\\
    &= \left[\frac{\theta}{2}- \frac{\sin{2\theta}}{4}\right]_0^{2\pi}
    &= \pi
\end{align*}

    \item $r= :\sin{u}$とおくと、$dr= \cos{u}du$。積分区間は\begin{tabular}[cl]$r$:&$0\to 1$\\ $\theta$:&$0\to \frac{\pi}{2}$\end{tabular}となる。
\begin{align*}
    \int_0^{\frac{\pi}{2}} \sin^2{u}\sqrt{1- \sin^2{2u}}\cdot \cos{u}du\\
    &= \int_0^{\frac{\pi}{2}} \sin^2{u}\cdot \cos^2{u}du\\
    &= \int_0^{\frac{\pi}{2}} \frac{\sin^2{2u}}{4}du\\
    &= \frac{1}{4}\int_0^{\frac{\pi}{2}}\frac{1- \cos{4u}}{2}du\\
    &= \frac{1}{8}\left[u- \frac{\sin{4u}}{4}\right]_0^{\frac{\pi}{2}}\\
    &= \frac{\pi}{16}
\end{align*}
    \item ヤコビアンの定義に従って計算すると
\begin{align*}
    \frac{\partial (x, y)}{\partial (r, \theta)}
    &= \begin{bmatrix}\partial x/\partial r& \partial x/\partial \theta\\ \partial y/\partial r& \partial y/\partial \theta\end{bmatrix}\\
    &= \begin{bmatrix}\cos{\theta}& -r\sin{\theta}\\ 2\ssin{\theta}& 2r\cos{\theta}\end{bmatrix}\\
    &= 2r
\end{align*}
    \item 変数変換することにより積分範囲は$0\le r\le 1, 0\le \theta\le 2\pi$となり、$dxdy= |2r|drd\theta= 2rdrd\theta$となる。
\begin{align*}
    \iint_D (x^2+ y^2)dxdy\\
    &= \int_{\theta= 0}^{2\pi}\int_{r= 0}^1 r^2(1+ 3\sin^2{theta})\cdot 2rdrd\theta\\
    &= \left(\int_{\theta= 0}^{2\pi} (1+ 3\sin^2{\theta})d\theta\right)\left(\int_{r= 0}^1 2r^3dr\right)\\
    &= \left([\theta]_{\theta= 0}^{2\pi}+ 3\pi\right)\left[\frac{r^4}{2}\right]_{r= 0}^1\\
    &= \frac{5\pi}{2}
\end{align*}
    \item 
\begin{align*}
    \frac{\partial (x, y, z)}{\partial (r, \theta, \phi)}
    &= 
    \begin{bmatrix}
        \sin{\theta}\cos{\phi}& r\cos{\theta}\cos{\phi}& -r\sin{\theta}\sin{\phi}\\
        \sin{\theta}\sin{\phi}& r\cos{\theta}\sin{\phi}& r\sin{\theta}\cos{\phi}\\
        \cos{\theta}& -r\sin{\theta}& 0
    \end{bmatrix}
    &= r^2\sin{\theta}
\end{align*}
    \item 変数変換を行うと、積分領域は$0\le r\le 1, 0\le \theta\le \pi, 0\le \phi\le 2\pi$となる。領域内で$\sin{\theta}\ge 0$であるため、体積要素は$dxdydz= r^2\sin{\theta}drd\theta d\phi$。
\begin{align*}
    \iiint_E 2\sqrt{1- x^2- y^2- z^2}dxdydz\\
    &= \int_{\phi= 0}^{2\pi}\int_{\theta= 0}^\pi\int_{r= 0}^1 2\sqrt{1- r^2}\cdot r^2\sin{\theta}drd\theta d\phi\\
    &= \left(\int_{\phi= 0}^{2\pi}d\phi\right)\left(\int_{\theta= 0}^\pi \sin{\theta}d\theta\right)\left(2\int_{r= 0}^1 r^2\sqrt{1- r^2}dr\right)\\
    &= 2\pi\times 2\times 2\times \frac{\pi}{16}\\
    &= \frac{\pi^2}{2}
\end{align*}
\end{enumerate}
\end{souanswer}
